\documentclass[12pt,a4paper]{article}
\usepackage[utf8]{inputenc}
\usepackage[T1]{fontenc}
\usepackage[brazilian]{babel}
\usepackage{geometry}
\usepackage{graphicx}
\usepackage{hyperref}
\usepackage{listings}
\usepackage{xcolor}
\usepackage{longtable}
\usepackage{booktabs}
\usepackage{array}
\usepackage{float}
\usepackage{caption}

\geometry{margin=2.5cm}

\definecolor{codegreen}{rgb}{0,0.6,0}
\definecolor{codegray}{rgb}{0.5,0.5,0.5}
\definecolor{codepurple}{rgb}{0.58,0,0.82}
\definecolor{backcolour}{rgb}{0.95,0.95,0.92}

\lstdefinestyle{sparqlstyle}{
    backgroundcolor=\color{backcolour},   
    commentstyle=\color{codegreen},
    keywordstyle=\color{blue},
    numberstyle=\tiny\color{codegray},
    stringstyle=\color{codepurple},
    basicstyle=\ttfamily\footnotesize,
    breakatwhitespace=false,         
    breaklines=true,                 
    captionpos=b,                    
    keepspaces=true,                 
    numbers=left,                    
    numbersep=5pt,                  
    showspaces=false,                
    showstringspaces=false,
    showtabs=false,                  
    tabsize=2,
    frame=single,
    morekeywords={SELECT, WHERE, PREFIX, ORDER, BY, FILTER, VALUES, DISTINCT, GROUP, LIMIT, DESC, ASC, UNION, MIN, MAX, AS}
}

\lstset{style=sparqlstyle}

\title{\textbf{MAC0444 - Projeto}\\
\large Relatório}
\author{Leonardo Heidi Almeida Murakami\\NUSP:11260186}
\date{\today}

\begin{document}

\maketitle

\tableofcontents
\newpage

% ============================================================================
\section{Introdução}
% ============================================================================

Este relatório descreve o desenvolvimento de uma ontologia para representação do domínio cinematográfico, incluindo conceitos como filmes, atores, atrizes, diretores e personagens. O projeto foi desenvolvido utilizando o Protégé como ferramenta principal, com referência à ontologia FOAF (\textit{Friend of a Friend}) como base para a modelagem de pessoas e suas relações.

A ontologia foi projetada para responder a dez questões de competência específicas, demonstrando sua capacidade de representar adequadamente o domínio do cinema e permitir consultas complexas sobre os dados.

% ============================================================================
\section{Descrição da Ontologia}
% ============================================================================

\subsection{URI Base}

A ontologia desenvolvida possui o seguinte URI base:

\begin{verbatim}
http://www.semanticweb.org/leonardomurakami/ontologies/2025/10/ontology-movies#
\end{verbatim}

\subsection{Importações}

A ontologia importa a ontologia FOAF (\textit{Friend of a Friend}), disponível em:

\begin{verbatim}
http://xmlns.com/foaf/0.1/
\end{verbatim}

ou em

\begin{verbatim}
https://web.archive.org/web/20211024051123/http://xmlns.com/foaf/spec/index.rdf    
\end{verbatim}

Esta importação é declarada como \texttt{owl:imports} no cabeçalho da ontologia, permitindo reutilizar conceitos e propriedades já estabelecidos para representação de pessoas e suas relações.

% ============================================================================
\section{Classes (Conceitos)}
% ============================================================================

A ontologia define cinco classes principais para representar as entidades do domínio cinematográfico:

\subsection{Movie}

\begin{itemize}
    \item \textbf{URI:} \texttt{:Movie}
    \item \textbf{Superclasse:} \texttt{owl:Thing}
    \item \textbf{Descrição:} Representa um filme. Cada instância possui título, ano de criação e duração, além de relações com diretores, atores, atrizes e personagens.
\end{itemize}

\subsection{Actor}

\begin{itemize}
    \item \textbf{URI:} \texttt{:Actor}
    \item \textbf{Superclasse:} \texttt{contact:Person} (via FOAF)
    \item \textbf{Descrição:} Representa um ator (pessoa do sexo masculino que atua em filmes). Possui primeiro nome, sobrenome e relações com os filmes em que atuou e os personagens que interpretou.
\end{itemize}

\subsection{Actress}

\begin{itemize}
    \item \textbf{URI:} \texttt{:Actress}
    \item \textbf{Superclasse:} \texttt{contact:Person} (via FOAF)
    \item \textbf{Descrição:} Representa uma atriz (pessoa do sexo feminino que atua em filmes). Possui as mesmas propriedades que Actor, mas permite distinguir gênero nas consultas.
\end{itemize}

\subsection{Director}

\begin{itemize}
    \item \textbf{URI:} \texttt{:Director}
    \item \textbf{Superclasse:} \texttt{contact:Person} (via FOAF)
    \item \textbf{Descrição:} Representa um diretor de filmes. Possui primeiro nome, sobrenome e relações com os filmes que dirigiu.
\end{itemize}

\subsection{Character}

\begin{itemize}
    \item \textbf{URI:} \texttt{:Character}
    \item \textbf{Superclasse:} \texttt{contact:Person} (via FOAF)
    \item \textbf{Descrição:} Representa um personagem de um filme. Cada personagem está associado a um ator ou atriz específico em um filme específico. Usamos a coluna id da planilha para evitar duplicação de personagens (ex: um doutor figurante, doutores de Doctor Who, etc.) 
\end{itemize}

\newpage
% ============================================================================
\section{Propriedades de Objeto (Object Properties)}
% ============================================================================

\subsection{Propriedades Derivadas de foaf:made}

Estas propriedades representam relações onde uma pessoa ``faz'' algo (atua, dirige):

\subsubsection{actsIn}
\begin{itemize}
    \item \textbf{URI:} \texttt{:actsIn}
    \item \textbf{Superpropriedade:} \texttt{foaf:made}
    \item \textbf{Domínio:} Actor $\cup$ Actress
    \item \textbf{Contradomínio:} Movie
    \item \textbf{Descrição:} Relaciona um ator/atriz ao filme em que ele/ela atua.
\end{itemize}

\subsubsection{actsAs}
\begin{itemize}
    \item \textbf{URI:} \texttt{:actsAs}
    \item \textbf{Superpropriedade:} \texttt{foaf:made}
    \item \textbf{Domínio:} Actor $\cup$ Actress
    \item \textbf{Contradomínio:} Character
    \item \textbf{Descrição:} Relaciona um ator/atriz ao personagem que ele/ela interpreta.
\end{itemize}

\subsubsection{directs}
\begin{itemize}
    \item \textbf{URI:} \texttt{:directs}
    \item \textbf{Superpropriedade:} \texttt{foaf:made}
    \item \textbf{Domínio:} Director
    \item \textbf{Contradomínio:} Movie
    \item \textbf{Descrição:} Relaciona um diretor ao filme que ele dirigiu.
\end{itemize}

\subsection{Propriedades Derivadas de foaf:maker}

Estas propriedades representam relações inversas, onde algo ``tem'' um criador:

\subsubsection{hasActor}
\begin{itemize}
    \item \textbf{URI:} \texttt{:hasActor}
    \item \textbf{Superpropriedade:} \texttt{foaf:maker}
    \item \textbf{Domínio:} Movie
    \item \textbf{Contradomínio:} Actor
    \item \textbf{Descrição:} Relaciona um filme a um ator que atua nele.
\end{itemize}

\subsubsection{hasActress}
\begin{itemize}
    \item \textbf{URI:} \texttt{:hasActress}
    \item \textbf{Superpropriedade:} \texttt{foaf:maker}
    \item \textbf{Domínio:} Movie
    \item \textbf{Contradomínio:} Actress
    \item \textbf{Descrição:} Relaciona um filme a uma atriz que atua nele.
\end{itemize}

\subsubsection{hasDirector}
\begin{itemize}
    \item \textbf{URI:} \texttt{:hasDirector}
    \item \textbf{Superpropriedade:} \texttt{foaf:maker}
    \item \textbf{Domínio:} Movie
    \item \textbf{Contradomínio:} Director
    \item \textbf{Descrição:} Relaciona um filme ao seu diretor.
\end{itemize}

\subsubsection{hasCharacter}
\begin{itemize}
    \item \textbf{URI:} \texttt{:hasCharacter}
    \item \textbf{Superpropriedade:} \texttt{foaf:maker}
    \item \textbf{Domínio:} Movie
    \item \textbf{Contradomínio:} Character
    \item \textbf{Descrição:} Relaciona um filme a um personagem que aparece nele.
\end{itemize}

\subsubsection{isPlayedBy}
\begin{itemize}
    \item \textbf{URI:} \texttt{:isPlayedBy}
    \item \textbf{Superpropriedade:} \texttt{foaf:maker}
    \item \textbf{Domínio:} Character
    \item \textbf{Contradomínio:} Actor $\cup$ Actress
    \item \textbf{Descrição:} Relaciona um personagem ao ator/atriz que o interpreta.
\end{itemize}

% ============================================================================
\section{Propriedades de Dados (Data Properties)}
% ============================================================================

\subsection{Propriedades da Ontologia FOAF Utilizadas}

\subsubsection{foaf:firstName}
\begin{itemize}
    \item \textbf{URI:} \texttt{foaf:firstName}
    \item \textbf{Tipo:} \texttt{xsd:string}
    \item \textbf{Descrição:} Primeiro nome de uma pessoa (ator, atriz ou diretor).
\end{itemize}

\subsubsection{foaf:familyName}
\begin{itemize}
    \item \textbf{URI:} \texttt{foaf:familyName}
    \item \textbf{Tipo:} \texttt{xsd:string}
    \item \textbf{Descrição:} Sobrenome (nome de família) de uma pessoa.
\end{itemize}

\subsubsection{foaf:title}
\begin{itemize}
    \item \textbf{URI:} \texttt{foaf:title}
    \item \textbf{Tipo:} \texttt{xsd:string}
    \item \textbf{Descrição:} Título do filme.
\end{itemize}

\subsection{Propriedades Definidas na Ontologia}

\subsubsection{creationYear}
\begin{itemize}
    \item \textbf{URI:} \texttt{:creationYear}
    \item \textbf{Tipo:} \texttt{xsd:integer}
    \item \textbf{Descrição:} Ano de lançamento do filme (pensei em criar como sub-propriedade de birthDay mas nao parecia fazer sentido).
\end{itemize}

\subsubsection{duration}
\begin{itemize}
    \item \textbf{URI:} \texttt{:duration}
    \item \textbf{Tipo:} \texttt{xsd:integer}
    \item \textbf{Descrição:} Duração do filme em minutos.
\end{itemize}

\subsubsection{birthDay}
\begin{itemize}
    \item \textbf{URI:} \texttt{:birthDay}
    \item \textbf{Superpropriedade:} \texttt{foaf:birthday}
    \item \textbf{Tipo:} \texttt{xsd:string}
    \item \textbf{Descrição:} Data de nascimento de uma pessoa.
\end{itemize}

% ============================================================================
\section{Relação com a Ontologia FOAF}
% ============================================================================

A ontologia desenvolvida segue rigorosamente os requisitos de integração com a FOAF, conforme especificado no enunciado do projeto:

\begin{table}[H]
\centering
\caption{Mapeamento entre elementos da ontologia e FOAF}
\begin{tabular}{|l|l|l|}
\hline
\textbf{Elemento FOAF} & \textbf{Uso na Ontologia} \\
\hline
\texttt{foaf:Person} & Superclasse de Actor, Actress, Director, Character  \\
\hline
\texttt{foaf:firstName} & Utilizado para primeiro nome \\
\hline
\texttt{foaf:familyName} & Utilizado para sobrenome  \\
\hline
\texttt{foaf:made} & Superpropriedade de actsIn, actsAs, directs \\
\hline
\texttt{foaf:maker} & Superpropriedade de hasActor, hasActress, etc. \\
\hline
\texttt{foaf:title} & Utilizado para título do filme \\
\hline
\end{tabular}
\end{table}

% ============================================================================
\section{Estatísticas da Ontologia}
% ============================================================================

A ontologia populada contém os seguintes números de indivíduos:

\begin{lstlisting}
SELECT ?Classe (COUNT(?individuo) AS ?Quantidade)
WHERE {
    ?individuo rdf:type ?Classe .
    FILTER(?Classe IN (:Movie, :Actor, :Actress, :Director, :Character))
}
GROUP BY ?Classe
ORDER BY DESC(?Quantidade)
\end{lstlisting}

\begin{table}[H]
\centering
\caption{Quantidade de indivíduos por classe}
\begin{tabular}{|l|r|}
\hline
\textbf{Classe} & \textbf{Quantidade} \\
\hline
Character & 2.789 \\
\hline
Actor & 1.693 \\
\hline
Actress & 828 \\
\hline
Director & 73 \\
\hline
Movie & 72 \\
\hline
\textbf{Total} & \textbf{5.455} \\
\hline
\end{tabular}
\end{table}

\newpage
% ============================================================================
\section{Consultas SPARQL e Resultados}
% ============================================================================

A seguir são apresentadas as dez consultas SPARQL desenvolvidas para responder às questões de competência, com seus respectivos parâmetros e resultados.

\subsection{Prefixos Utilizados}

\begin{lstlisting}
PREFIX : <http://www.semanticweb.org/leonardomurakami/
         ontologies/2025/10/ontology-movies#>
PREFIX rdf: <http://www.w3.org/1999/02/22-rdf-syntax-ns#>
PREFIX rdfs: <http://www.w3.org/2000/01/rdf-schema#>
PREFIX foaf: <http://xmlns.com/foaf/0.1/>
PREFIX xsd: <http://www.w3.org/2001/XMLSchema#>
\end{lstlisting}


% ----------------------------------------------------------------------------
\subsection{Query 1: Filmes dirigidos por um diretor}
% ----------------------------------------------------------------------------

\textbf{Questão:} Quais os títulos dos filmes que foram dirigidos pelo diretor D, em ordem lexicográfica?

\textbf{Parâmetro:} \texttt{D = :joel\_coen}

\begin{lstlisting}
SELECT ?TituloDoFilme
WHERE {
    VALUES ?D { :joel_coen }
    
    ?filme rdf:type :Movie .
    ?filme :hasDirector ?D .
    ?filme foaf:title ?TituloDoFilme .
}
ORDER BY ?TituloDoFilme
\end{lstlisting}

\textbf{Resultados:}

\begin{table}[H]
\centering
\begin{tabular}{|l|}
\hline
\textbf{TituloDoFilme} \\
\hline
Barton Fink \\
\hline
Blood Simple \\
\hline
Fargo \\
\hline
Intolerable Cruelty \\
\hline
Miller's Crossing \\
\hline
No Country for Old Men \\
\hline
O Brother, Where Art Thou? \\
\hline
Paris, je t'aime \\
\hline
Raising Arizona \\
\hline
The Big Lebowski \\
\hline
The Hudsucker Proxy \\
\hline
The Ladykillers \\
\hline
The Man Who Wasn't There \\
\hline
\end{tabular}
\end{table}

\newpage
% ----------------------------------------------------------------------------
\subsection{Query 2: Atores e personagens de um filme}
% ----------------------------------------------------------------------------

\textbf{Questão:} Quais os primeiros nomes e últimos nomes dos atores que participaram do filme de título Ft, em ordem lexicográfica de nome e sobrenome, com seus respectivos personagens?

\textbf{Parâmetro:} \texttt{Ft = "The Godfather"}

\begin{lstlisting}
SELECT ?PrimeiroNome ?UltimoNome ?Personagem
WHERE {
    VALUES ?Ft { "The Godfather" }
    
    ?filme rdf:type :Movie .
    ?filme foaf:title ?Ft .
    
    { ?filme :hasActor ?pessoa . }
    UNION
    { ?filme :hasActress ?pessoa . }
    
    ?pessoa foaf:firstName ?PrimeiroNome .
    ?pessoa foaf:familyName ?UltimoNome .
    
    ?pessoa :actsAs ?Personagem .
    ?filme :hasCharacter ?Personagem .
}
ORDER BY ?PrimeiroNome ?UltimoNome
\end{lstlisting}

\textbf{Resultados:}

\begin{longtable}{|l|l|p{6cm}|}
\hline
\textbf{PrimeiroNome} & \textbf{UltimoNome} & \textbf{Personagem} \\
\hline
\endfirsthead
\hline
\textbf{PrimeiroNome} & \textbf{UltimoNome} & \textbf{Personagem} \\
\hline
\endhead
Abe & Vigoda & abe\_vigodathe\_godfather salvadore\_sally\_tessio \\
\hline
Al & Lettieri & al\_lettierithe\_godfather virgil\_sollozzo \\
\hline
Al & Martino & al\_martinothe\_godfather johnny\_fontane \\
\hline
Al & Pacino & al\_pacinothe\_godfather michael\_corleone \\
\hline
Alex & Rocco & alex\_roccothe\_godfather moe\_greene \\
\hline
Angelo & Infanti & angelo\_infantithe\_godfather fabrizio \\
\hline
Anthony & Gounaris & anthony\_gounaristhe\_godfather anthony\_vito\_corleone \\
\hline
Ardell & Sheridan & ardell\_sheridanthe\_godfather mrs\_clemenza \\
\hline
Burt & Richards & burt\_richardsthe\_godfather floral\_designer \\
\hline
Carmine & Coppola & carmine\_coppolathe\_godfather piano\_player\_in\_montage\_scene \\
\hline
Corrado & Gaipa & corrado\_gaipathe\_godfather don\_tommasino \\
\hline
Diane & Keaton & diane\_keatonthe\_godfather kay\_adams \\
\hline
Ed & Vantura & ed\_vanturathe\_godfather wedding\_guest \\
\hline
Father & Medeglia & father\_joseph\_medegliathe\_godfather priest\_at\_baptism \\
\hline
Filomena & Spagnuolo & filomena\_spagnuolothe\_godfather extra\_at\_wedding\_scene \\
\hline
Franco & Citti & franco\_cittithe\_godfather calo \\
\hline
Frank & Macetta & frank\_macettathe\_godfather \\
\hline
Frank & Sivero & frank\_siverothe\_godfather extra \\
\hline
Gabriele & Torrei & gabriele\_torreithe\_godfather enzo\_robutti\_the\_baker \\
\hline
Gian & Coppola & gian\_carlo\_coppolathe\_godfather baptism\_observer \\
\hline
Gianni & Russo & gianni\_russothe\_godfather carlo\_rizzi \\
\hline
James & Caan & james\_caanthe\_godfather santino\_sonny\_corleone \\
\hline
Jeannie & Linero & jeannie\_linerothe\_godfather lucy\_mancini \\
\hline
Joe & Grippo & joe\_lo\_grippothe\_godfather sonny\_s\_bodyguard \\
\hline
Joe & Spinell & joe\_spinellthe\_godfather willie\_cicci \\
\hline
John & Cazale & john\_cazalethe\_godfather fredo \\
\hline
John & Marley & john\_marleythe\_godfather jack\_woltz \\
\hline
John & Martino & john\_martinothe\_godfather paulie\_gatto \\
\hline
Julie & Gregg & julie\_greggthe\_godfather sandra\_corleone \\
\hline
Lenny & Montana & lenny\_montanathe\_godfather luca\_brasi \\
\hline
Lou & Jr. & lou\_martini\_jrthe\_godfather boy\_at\_wedding \\
\hline
Louis & Guss & louis\_gussthe\_godfather don\_zaluchi \\
\hline
Marlon & Brando & marlon\_brandothe\_godfather don\_vito\_corleone \\
\hline
Matthew & Vlahakis & matthew\_vlahakisthe\_godfather clemenza\_s\_son \\
\hline
Morgana & King & morgana\_kingthe\_godfather mama\_corleone \\
\hline
Nick & Vallelonga & nick\_vallelongathe\_godfather wedding\_party\_guest \\
\hline
Randy & Jurgensen & randy\_jurgensenthe\_godfather sonny\_s\_killer\_1 \\
\hline
Richard & Bright & richard\_brightthe\_godfather al\_neri \\
\hline
Richard & Castellano & richard\_s\_castellanothe\_godfather pete\_clemenza \\
\hline
Richard & Conte & richard\_contethe\_godfather emilio\_barzini \\
\hline
Rick & Petrucelli & rick\_petrucellithe\_godfather man\_in\_passenger\_seat \\
\hline
Robert & Duvall & robert\_duvallthe\_godfather tom\_hagen \\
\hline
Ron & Gilbert & ron\_gilbertthe\_godfather usher\_in\_bridal\_party \\
\hline
Rudy & Bond & rudy\_bondthe\_godfather ottilio\_cuneo \\
\hline
Sal & Richards & sal\_richardsthe\_godfather drunk \\
\hline
Salvatore & Corsitto & salvatore\_corsittothe\_godfather amerigo\_bonasera \\
\hline
Saro & Urz & saro\_urzthe\_godfather vitelli \\
\hline
Simonetta & Stefanelli & simonetta\_stefanellithe\_godfather apollonia\_vitelli\_corleone \\
\hline
Sofia & Coppola & sofia\_coppolathe\_godfather michael\_francis\_rizzi \\
\hline
Sonny & Grosso & sonny\_grossothe\_godfather cop\_with\_capt\_mccluskey \\
\hline
Sterling & Hayden & sterling\_haydenthe\_godfather capt\_mark\_mccluskey \\
\hline
Talia & Shire & talia\_shirethe\_godfather connie \\
\hline
Tere & Livrano & tere\_livranothe\_godfather theresa\_hagen \\
\hline
Tom & Rosqui & tom\_rosquithe\_godfather rocco\_lampone \\
\hline
Tony & Giorgio & tony\_giorgiothe\_godfather bruno\_tattaglia \\
\hline
Tony & Lip & tony\_lipthe\_godfather wedding\_guest \\
\hline
Victor & Rendina & victor\_rendinathe\_godfather philip\_tattaglia \\
\hline
Vito & Scotti & vito\_scottithe\_godfather nazorine \\
\hline
\end{longtable}
\newpage
% ----------------------------------------------------------------------------
\subsection{Query 3: Filmes em que dois atores atuaram juntos}
% ----------------------------------------------------------------------------

\textbf{Questão:} Em quais filmes os atores X e Y atuaram juntos, com os respectivos diretores e anos de lançamento, do mais novo para o mais antigo?

\textbf{Parâmetros:} \texttt{X = :al\_pacino}, \texttt{Y = :diane\_keaton}

\begin{lstlisting}
SELECT ?Filme ?Diretor ?Ano
WHERE {
    VALUES ?X { :al_pacino }
    VALUES ?Y { :diane_keaton }
    
    ?filme rdf:type :Movie .
    ?filme foaf:title ?Filme .
    ?filme :creationYear ?Ano .
    ?filme :hasDirector ?Diretor .
    
    { ?filme :hasActor ?X . } UNION { ?filme :hasActress ?X . }
    { ?filme :hasActor ?Y . } UNION { ?filme :hasActress ?Y . }
}
ORDER BY DESC(?Ano)
\end{lstlisting}

\textbf{Resultados:}

\begin{table}[H]
\centering
\begin{tabular}{|l|l|r|}
\hline
\textbf{Filme} & \textbf{Diretor} & \textbf{Ano} \\
\hline
The Godfather: Part III & :francis\_ford\_coppola & 1990 \\
\hline
The Godfather: Part II & :francis\_ford\_coppola & 1974 \\
\hline
The Godfather & :francis\_ford\_coppola & 1972 \\
\hline
\end{tabular}
\end{table}
\newpage
% ----------------------------------------------------------------------------
\subsection{Query 4: Filmes do diretor de F com atores X ou Y}
% ----------------------------------------------------------------------------

\textbf{Questão:} Quais filmes do diretor do filme F possuem X ou Y como atores, com as respectivas durações em ordem crescente?

\textbf{Parâmetros:} \texttt{F = "The Godfather"}, \texttt{X = :al\_pacino}, \texttt{Y = :diane\_keaton}

\begin{lstlisting}
SELECT DISTINCT ?Filme ?Duracao
WHERE {
    VALUES ?F { "The Godfather" }
    VALUES ?X { :al_pacino }
    VALUES ?Y { :diane_keaton }
    
    ?filmeF rdf:type :Movie .
    ?filmeF foaf:title ?F .
    ?filmeF :hasDirector ?diretor .
    
    ?filme :hasDirector ?diretor .
    ?filme foaf:title ?Filme .
    ?filme :duration ?Duracao .
    
    {
        { ?filme :hasActor ?X . } UNION { ?filme :hasActress ?X . }
    }
    UNION
    {
        { ?filme :hasActor ?Y . } UNION { ?filme :hasActress ?Y . }
    }
}
ORDER BY ?Duracao
\end{lstlisting}

\textbf{Resultados:}

\begin{table}[H]
\centering
\begin{tabular}{|l|r|}
\hline
\textbf{Filme} & \textbf{Duracao} \\
\hline
The Godfather: Part III & 162 \\
\hline
The Godfather & 175 \\
\hline
The Godfather: Part II & 202 \\
\hline
\end{tabular}
\end{table}
\newpage
% ----------------------------------------------------------------------------
\subsection{Query 5: Pessoas que atuaram e dirigiram}
% ----------------------------------------------------------------------------

\textbf{Questão:} Quais pessoas atuaram em um filme e dirigiram algum filme (não necessariamente o mesmo filme, mas obrigatoriamente a mesma pessoa)?

\textbf{Parâmetros:} Nenhum (busca todas as pessoas).

\begin{lstlisting}
SELECT DISTINCT ?Pessoa ?FilmeAt ?FilmeDir
WHERE {
    { ?filmeAtuado :hasActor ?Pessoa . }
    UNION
    { ?filmeAtuado :hasActress ?Pessoa . }
    
    ?filmeDirigido :hasDirector ?Pessoa .
    
    ?filmeAtuado foaf:title ?FilmeAt .
    ?filmeDirigido foaf:title ?FilmeDir .
}
ORDER BY ?Pessoa ?FilmeAt ?FilmeDir
\end{lstlisting}

\textbf{Resultados (amostra representativa):}

\begin{longtable}{|l|p{4.5cm}|p{4.5cm}|}
\hline
\textbf{Pessoa} & \textbf{FilmeAt} & \textbf{FilmeDir} \\
\hline
\endfirsthead
\hline
\textbf{Pessoa} & \textbf{FilmeAt} & \textbf{FilmeDir} \\
\hline
\endhead
:denzel\_washington & The Great Debaters & The Great Debaters \\
\hline
:eleanor\_coppola & Hearts of Darkness: A Filmmaker's Apocalypse & Hearts of Darkness: A Filmmaker's Apocalypse \\
\hline
:eric\_schaeffer & Fall & Fall \\
\hline
:eric\_schaeffer & Fall & If Lucy Fell \\
\hline
:eric\_schaeffer & If Lucy Fell & Fall \\
\hline
:eric\_schaeffer & If Lucy Fell & If Lucy Fell \\
\hline
:ethan\_coen & Down from the Mountain & Barton Fink \\
\hline
:ethan\_coen & Down from the Mountain & Blood Simple \\
\hline
:ethan\_coen & Down from the Mountain & Fargo \\
\hline
:francis\_ford\_coppola & Hearts of Darkness: A Filmmaker's Apocalypse & The Godfather \\
\hline
:francis\_ford\_coppola & Hearts of Darkness: A Filmmaker's Apocalypse & The Godfather: Part II \\
\hline
:francis\_ford\_coppola & Hearts of Darkness: A Filmmaker's Apocalypse & The Godfather: Part III \\
\hline
:george\_lucas & Hearts of Darkness: A Filmmaker's Apocalypse & Star Wars: Episode I \\
\hline
:harold\_ramis & Ghostbusters & Groundhog Day \\
\hline
:harold\_ramis & Groundhog Day & Groundhog Day \\
\hline
:ivan\_reitman & Ghostbusters & Ghostbusters \\
\hline
:joel\_coen & Crimewave & Blood Simple \\
\hline
:joel\_coen & Crimewave & Fargo \\
\hline
:joel\_coen & Down from the Mountain & The Big Lebowski \\
\hline
:robert\_redford & The Horse Whisperer & The Horse Whisperer \\
\hline
:roman\_coppola & Hearts of Darkness: A Filmmaker's Apocalypse & CQ \\
\hline
:sam\_raimi & Miller's Crossing & Crimewave \\
\hline
:sam\_raimi & The Hudsucker Proxy & Crimewave \\
\hline
:sofia\_coppola & The Godfather & Lost in Translation \\
\hline
:sofia\_coppola & The Godfather & Marie Antoinette \\
\hline
:sofia\_coppola & The Godfather: Part III & The Virgin Suicides \\
\hline
:spike\_jonze & Torrance Rises & Her \\
\hline
:stephen\_hillenburg & The SpongeBob SquarePants Movie & The SpongeBob SquarePants Movie \\
\hline
:woody\_allen & Scoop & Match Point \\
\hline
:woody\_allen & Scoop & Scoop \\
\hline
\end{longtable}

\textit{Nota: A tabela completa contém 138 linhas com todas as combinações de filmes atuados e dirigidos por cada pessoa.}
\newpage
% ----------------------------------------------------------------------------
\subsection{Query 6: Diretores em intervalo de anos com atores X e Y}
% ----------------------------------------------------------------------------

\textbf{Questão:} Quais diretores dirigiram algum filme em um ano entre os anos N1 e N2, em que os atores X e Y aparecem, do mais antigo para o mais novo?

\textbf{Parâmetros:} \texttt{N1 = 1940}, \texttt{N2 = 2000}, \texttt{X = :al\_pacino}, \texttt{Y = :al\_martino}

\begin{lstlisting}
SELECT ?Ano ?Diretor ?Filme
WHERE {
    VALUES ?N1 { 1940 }
    VALUES ?N2 { 2000 }
    VALUES ?X { :al_pacino }
    VALUES ?Y { :al_martino }
    
    ?filme rdf:type :Movie .
    ?filme foaf:title ?Filme .
    ?filme :creationYear ?Ano .
    ?filme :hasDirector ?Diretor .
    
    FILTER(?Ano >= ?N1 && ?Ano <= ?N2)
    
    { ?filme :hasActor ?X . } UNION { ?filme :hasActress ?X . }
    { ?filme :hasActor ?Y . } UNION { ?filme :hasActress ?Y . }
}
ORDER BY ?Ano
\end{lstlisting}

\textbf{Resultados:}

\begin{table}[H]
\centering
\begin{tabular}{|r|l|l|}
\hline
\textbf{Ano} & \textbf{Diretor} & \textbf{Filme} \\
\hline
1972 & :francis\_ford\_coppola & The Godfather \\
\hline
1990 & :francis\_ford\_coppola & The Godfather Part III \\
\hline
\end{tabular}
\end{table}
\newpage
% ----------------------------------------------------------------------------
\subsection{Query 7: Pares atriz-ator em filmes de duração específica}
% ----------------------------------------------------------------------------

\textbf{Questão:} Quais pares formados por uma atriz e um ator atuaram juntos em algum filme de duração entre M1 e M2?

\textbf{Parâmetros:} \texttt{M1 = 140}, \texttt{M2 = 150}

\begin{lstlisting}
SELECT DISTINCT ?Atriz ?Ator ?Filme ?Duracao
WHERE {
    VALUES ?M1 { 140 }
    VALUES ?M2 { 150 }
    
    ?Filme rdf:type :Movie .
    ?Filme :duration ?Duracao .
    ?Filme :hasActress ?Atriz .
    ?Filme :hasActor ?Ator .
    
    FILTER(?Duracao >= ?M1 && ?Duracao <= ?M2)
}
ORDER BY ?Atriz ?Ator
\end{lstlisting}

\textbf{Resultados:}

\begin{table}[H]
\centering
\footnotesize
\begin{tabular}{|l|l|l|r|}
\hline
\textbf{Atriz} & \textbf{Ator} & \textbf{Filme} & \textbf{Duracao} \\
\hline
:audrey\_tautou & :alfred\_molina & :the\_da\_vinci\_code & 149 \\
\hline
:audrey\_tautou & :ian\_mckellen & :the\_da\_vinci\_code & 149 \\
\hline
:audrey\_tautou & :jean\_reno & :the\_da\_vinci\_code & 149 \\
\hline
:audrey\_tautou & :jean\_yves\_berteloot & :the\_da\_vinci\_code & 149 \\
\hline
:audrey\_tautou & :jurgen\_prochnow & :the\_da\_vinci\_code & 149 \\
\hline
:audrey\_tautou & :paul\_bettany & :the\_da\_vinci\_code & 149 \\
\hline
:audrey\_tautou & :tom\_hanks & :the\_da\_vinci\_code & 149 \\
\hline
\end{tabular}
\end{table}
\newpage
% ----------------------------------------------------------------------------
\subsection{Query 8: Diretor que mais dirigiu filmes com ator específico}
% ----------------------------------------------------------------------------

\textbf{Questão:} Qual o diretor que mais dirigiu filmes em que aparece o ator de primeiro nome Xp e último nome Xu?

\textbf{Parâmetros:} \texttt{Xp = "Al"}, \texttt{Xu = "Pacino"}

\begin{lstlisting}
SELECT ?Diretor (COUNT(?filme) AS ?NumFilmes)
WHERE {
    VALUES ?Xp { "Al" }
    VALUES ?Xu { "Pacino" }
    
    ?ator foaf:firstName ?Xp .
    ?ator foaf:familyName ?Xu .
    
    { ?filme :hasActor ?ator . } 
    UNION 
    { ?filme :hasActress ?ator . }
    
    ?filme :hasDirector ?Diretor .
}
GROUP BY ?Diretor
ORDER BY DESC(?NumFilmes)
LIMIT 1
\end{lstlisting}

\textbf{Resultado:}

\begin{table}[H]
\centering
\begin{tabular}{|l|r|}
\hline
\textbf{Diretor} & \textbf{numFilmes} \\
\hline
:francis\_ford\_coppola & 3 \\
\hline
\end{tabular}
\end{table}
\newpage
% ----------------------------------------------------------------------------
\subsection{Query 9: Filme mais antigo de um ator}
% ----------------------------------------------------------------------------

\textbf{Questão:} Qual o filme mais antigo em que o ator X atuou, em que ano foi lançado e qual era seu personagem? Havendo empate, devolver todos do mesmo ano.

\textbf{Parâmetro:} \texttt{X = :joel\_coen}

\begin{lstlisting}
SELECT ?Filme ?Ano ?Personagem
WHERE {
    VALUES ?X { :joel_coen }
    
    ?Filme rdf:type :Movie .
    ?Filme foaf:title ?tituloFilme .
    ?Filme :creationYear ?Ano .
    
    { ?Filme :hasActor ?X . } UNION { ?Filme :hasActress ?X . }
    
    ?X :actsAs ?Personagem .
    ?Filme :hasCharacter ?Personagem .
    
    FILTER NOT EXISTS {
        ?outroFilme :creationYear ?outroAno .
        { ?outroFilme :hasActor ?X . } UNION { ?outroFilme :hasActress ?X . }
        FILTER(?outroAno < ?Ano)
    }
}
\end{lstlisting}

\textbf{Resultado:}

\begin{table}[H]
\centering
\begin{tabular}{|l|r|l|}
\hline
\textbf{Filme} & \textbf{Ano} & \textbf{Personagem} \\
\hline
:spies\_like\_us & 1985 & :joel\_coenspies\_like\_us drive\_in\_security \\
\hline
:crimewave & 1985 & :joel\_coencrimewave reporter\_at\_execution \\
\hline
\end{tabular}
\end{table}
\newpage
% ----------------------------------------------------------------------------
\subsection{Query 10: Filme mais longo de um diretor}
% ----------------------------------------------------------------------------

\textbf{Questão:} Qual o filme mais longo dirigido por D, e qual a sua duração? Havendo empate, devolver todos da mesma duração.

\textbf{Parâmetro:} \texttt{D = :francis\_ford\_coppola}

\begin{lstlisting}
SELECT ?Filme ?Duracao
WHERE {
    VALUES ?D { :francis_ford_coppola }
    
    ?Filme rdf:type :Movie .
    ?Filme :hasDirector ?D .
    ?Filme :duration ?Duracao .
    
    {
        SELECT (MAX(?duracaoMax) AS ?duracaoMaxima)
        WHERE {
            ?filmeInner :hasDirector ?D .
            ?filmeInner :duration ?duracaoMax .
        }
    }
    
    FILTER(?Duracao = ?duracaoMaxima)
}
\end{lstlisting}

\textbf{Resultado:}

\begin{table}[H]
\centering
\begin{tabular}{|l|r|}
\hline
\textbf{Filme} & \textbf{Duracao} \\
\hline
:the\_godfather\_part\_ii & 202 \\
\hline
\end{tabular}
\end{table}
\end{document}